\documentclass[utf8,a4paper,fontset=none]{ctexart}

\usepackage{anyfontsize}
\usepackage{amsmath,amssymb,mathrsfs,wasysym}
\usepackage{autobreak}
\allowdisplaybreaks
\usepackage[T1]{fontenc}
\usepackage{textcomp}
\usepackage[libertine,vvarbb]{newtxmath}
\usepackage[scr=rsfso,frak=euler,bb=ams]{mathalfa}
\usepackage{bm}
\usepackage{indentfirst}
\setlength{\parindent}{0em}
\usepackage{parskip}
\usepackage{geometry}
\geometry{top=3cm,bottom=3cm,left=2cm,right=2cm}
\usepackage{fontspec}
\setmainfont{Linux Libertine O}
\setsansfont{Linux Biolinum O}
\setmonofont{JetBrains Mono}
\setCJKmainfont{SimSun}[BoldFont=Noto Sans CJK SC Medium]

\begin{document}

\title{\textbf{原子物理学作业}}
\author{1900011604\ \ 张植竣\ \ 北京大学}
\date{2021年5月20日}

\maketitle

\subsection*{补充题}

\begin{itemize}

\item[1.] 由于不存在稳定的质量数为5或8的原子核，早期宇宙无法从$^4\mathrm{He}$或$^7\mathrm{Li}$合成$^{12}\mathrm{C}$，而三个$^4\mathrm{He}$碰撞形成$^{12}\mathrm{C}$由于温度和压力太低而很难发生。

\item[2.] 微量元素主要由$r$过程（双中子星并合、超新星爆发）和$s$过程（处于渐近巨星支的恒星）产生，例如I主要在双中子星并合（$r$过程）中合成，而Zn和Se主要形成于超新星爆发（$r$过程）中。

\item[3.(1)] 不参与电磁相互作用的有中子n和中微子$\nu$。

\item[3.(2)] 不参与强相互作用的有电子e、光子$\gamma$和中微子$\nu$。

\item[4.] 因为常见的强子奇异数均为0，如$\pi^+$介子，$\pi^0$介子，质子和中子，而奇异子非常不稳定，寿命约在$10^{-10}\,\mathrm{s}$的量级。在高能条件下发现奇异粒子之前，常规的能量较低的日常物质只有两味。

\item[5.]

\end{itemize}

\end{document}